\documentclass[11pt,a4paper]{article}
\usepackage[T1]{fontenc}
\usepackage[utf8]{inputenc}
\usepackage{lmodern,microtype}
\usepackage[margin=27mm]{geometry}
\usepackage{amsmath,amssymb,amsthm,mathtools}
\usepackage{booktabs,longtable,array,enumitem}
\usepackage{tikz-cd}
\usepackage{xcolor}
\usepackage[colorlinks=true,linkcolor=blue!55!black,citecolor=blue!55!black,urlcolor=blue!55!black]{hyperref}
\usepackage{bookmark}
\hypersetup{pdftitle={A first introduction to Deligne--Mumford stacks},pdfauthor={Notes prepared for the AlgebraicCutting project}}
\setlength{\parskip}{4pt}
\setlength{\parindent}{0pt}
\setlength{\emergencystretch}{2em}
\setlist{itemsep=3pt,topsep=5pt}
\numberwithin{equation}{section}
\newtheorem{proposition}{Proposition}[section]
\newtheorem{theorem}[proposition]{Theorem}
\newtheorem{definition}[proposition]{Definition}
\newtheorem{lemma}[proposition]{Lemma}
\theoremstyle{definition}
\newtheorem{example}[proposition]{Example}
\newtheorem{exercise}{Exercise}
\newtheorem{remark}[proposition]{Remark}
\newcommand{\A}{\mathbb A}
\newcommand{\PP}{\mathbb P}
\newcommand{\Gm}{\mathbb G_m}
\newcommand{\OO}{\mathcal O}
\newcommand{\cX}{\mathcal X}
\newcommand{\cY}{\mathcal Y}
\newcommand{\cZ}{\mathcal Z}
\newcommand{\cD}{\mathcal D}
\newcommand{\cL}{\mathcal L}
\newcommand{\cM}{\mathcal M}
\newcommand{\cP}{\mathcal P}
\newcommand{\Spec}{\operatorname{Spec}}
\newcommand{\Proj}{\operatorname{Proj}}
\newcommand{\Hom}{\operatorname{Hom}}
\newcommand{\Aut}{\operatorname{Aut}}
\newcommand{\Isom}{\operatorname{Isom}}
\newcommand{\Stab}{\operatorname{Stab}}
\newcommand{\Pic}{\operatorname{Pic}}
\newcommand{\sst}{\mathrm{ss}}
\newcommand{\id}{\mathrm{id}}
\newcommand{\one}{\mathbf 1}
\newcommand{\normal}{\mathcal N}
\newcommand{\charline}[1]{k_{#1}}
\newcommand{\reading}[1]{\par\smallskip\textit{Reading checkpoint.} #1\par}
\newcommand{\source}[1]{\par\smallskip{\small\textit{Reference.} #1}\par}
\title{A first introduction to Deligne--Mumford stacks\\[5pt]
\large Quotients, families, line bundles, and algebraic cutting}
\author{Notes prepared for the AlgebraicCutting project}
\date{13 November 2026}

\begin{document}
\maketitle
\begin{abstract}
These notes start from schemes, gluing, and line bundles, at the level of
Hartshorne, Chapters I--II. They explain why quotient spaces can lose geometric
information, how groupoids and descent retain it, and what the definition of a
Deligne--Mumford stack says. A small collection of examples is developed repeatedly:
the classifying stack $B\mu_n$, the quotient $[\A^1/\mu_n]$, a quadratic cone,
and the weighted projective line $\PP(n,1)$. The last example is also an explicit
algebraic cut. It makes the distinction between a stack line bundle, its
pushforward to the coarse space, and a line bundle descended to that space visible
in elementary calculations. Definitions and worked arguments are separated from
general theorems used without proof.
\end{abstract}
\tableofcontents
\newpage

\section{How to read these notes}
\label{sec:reading}

We work over an algebraically closed field $k$ of characteristic zero; taking
$k=\mathbb C$ throughout is entirely sufficient. Schemes and group actions in the
examples are of finite type. The group $\Gm$ is $\Spec k[t,t^{-1}]$, and
$\mu_n\subset\Gm$ is the group of $n$th roots of unity. Here $n\geq2$ unless
otherwise stated.

You should be comfortable with affine schemes, fibre products, sheaves, line
bundles, divisors, and $\Proj$. No knowledge of Hartshorne, Chapter III, is assumed.
Some ingredients that are not developed systematically in Chapters I--II---
\emph{étale morphisms}, \emph{torsors}, and \emph{algebraic spaces}---are introduced
when they are first needed. Their general theory is unnecessary for the worked
examples. Higher categories, derived geometry, intersection theory on stacks, and
the construction of moduli of curves are not prerequisites.

The main route is
\[
\begin{gathered}
\text{quotient rings}\longrightarrow\text{objects and their isomorphisms}
\longrightarrow\text{families and gluing}\\
\longrightarrow\text{quotient stacks}\longrightarrow\text{DM stacks}
\longrightarrow\text{line bundles and cuts}.
\end{gathered}
\]
Sections~\ref{sec:quotients}--\ref{sec:dm} give the foundations.
Sections~\ref{sec:coarse}--\ref{sec:roots} work out the geometry and sheaves.
Sections~\ref{sec:cutting}--\ref{sec:toycut} connect them to semiprojective
algebraic cutting. The exercises in Section~\ref{sec:exercises} have solutions.

The notes use calligraphic letters for stacks and roman letters for their coarse
spaces. The notation $[X/G]$ always denotes a \emph{quotient stack}; an affine
invariant quotient is written $\Spec A^G$. These are different constructions.

\paragraph{A useful question to keep asking.}
When replacing an object by its isomorphism class, what have we forgotten about
its automorphisms, its families, or its line bundles? Each example below supplies
a concrete answer.

\section{What an invariant quotient forgets}
\label{sec:quotients}

\subsection{A quotient whose coarse space is smooth}
Let $\mu_n$ act on $\A^1_z$ by $z\mapsto\zeta z$. A polynomial is invariant
exactly when its monomials have degrees divisible by $n$, so
\[
 k[z]^{\mu_n}=k[z^n]=k[u],\qquad u=z^n.
\]
The usual affine quotient is another affine line,
\[
 q:\A^1_z\longrightarrow\A^1_u,\qquad z\longmapsto z^n.
\]
Over a nonzero geometric point $u$, the $n$ roots form one orbit and the action
on that orbit is free. At $z=0$, the orbit contains one point, but every
$\zeta\in\mu_n$ fixes it. Thus
\[
 \Stab(z)=\begin{cases}\{1\},&z\ne0,\\ \mu_n,&z=0.\end{cases}
\]
The coordinate ring of the quotient records invariant functions. It does not,
by itself, record this jump in stabilisers. In particular, smoothness of the
quotient $\A^1_u$ does not imply that the action was free.

One might try to decorate $u=0$ with the label $\mu_n$. That is a useful first
picture, but it is incomplete: we also need to say how the group acts on nearby
directions and on fibres of sheaves, and how such objects vary in families.

\subsection{A quotient whose coarse space is singular}
Let $\mu_2$ act on $\A^2_{x,y}$ by
\[
 (x,y)\longmapsto(-x,-y).
\]
An invariant monomial has even total degree. Hence
\[
 k[x,y]^{\mu_2}=k[a,b,c]/(ac-b^2),\qquad
 a=x^2,\quad b=xy,\quad c=y^2.
\]
For example, every even monomial is a product of $x^2,xy,y^2$, and the displayed
relation gives the familiar quadratic cone. The derivatives of $ac-b^2$ are
$c,-2b,a$, all zero at the vertex. The quotient is singular there, although
$\A^2$ is smooth.

We will obtain a smooth stack $[\A^2/\mu_2]$ with this cone as coarse space.
Its origin retains the sign action on the two-dimensional tangent space. The
previous example will give a smooth stack with a smooth coarse space. Taken
together, the examples show that the singularities of the coarse space do not
provide a complete test for the presence of stack structure.

\subsection{Retain arrows as well as orbits}
For a group $G$ acting on a set $S$, form a category with objects $x\in S$ and
arrows
\[
 \Hom(x,x')=\{g\in G:g\cdot x=x'\}.
\]
Composition is multiplication in $G$. Every arrow is invertible; a category with
this property is a \emph{groupoid}. The isomorphism classes are the orbits, while
$\Aut(x)=\Stab_G(x)$. Thus a groupoid keeps both the quotient points and their
symmetries.

There is an important scheme-level distinction. For the sign action on $\A^1$,
the action groupoid has arrow scheme $\mu_2\times\A^1$, which is two disjoint
affine lines. The relation obtained by comparing images in the coarse quotient
has scheme
\[
 \A^1\times_{\A^1_u}\A^1
 =\Spec k[z,z']/(z^2-z'^2).
\]
This is two lines meeting at the origin. The action groupoid has two distinct
arrows from $0$ to itself: the identity and the sign. The coarse relation has
lost their distinction. Replacing a groupoid by a relation is already a loss of
scheme-theoretic information.

\reading{Compute the stabilisers in both examples. Why could neither coarse
coordinate ring tell you which character occurs on an extra line over the origin?}

\section{From points to families: why a stack is groupoid-valued}
\label{sec:families}

\subsection{The functor-of-points viewpoint}
A scheme $X$ determines a contravariant functor
\[
 T\longmapsto X(T)=\Hom(T,X).
\]
A morphism $T'\to T$ pulls a $T$-point back to a $T'$-point. The Yoneda principle
says that the functor, with its behaviour on morphisms, determines $X$. One must
test on general schemes $T$, including nonreduced schemes, rather than just on
$\Spec k$.

For a moduli problem, a $T$-point is usually a \emph{family over $T$}. Two
families can be isomorphic in more than one way. The natural value of the moduli
problem is therefore a groupoid, not merely a set of isomorphism classes.

For example, let $\mathcal B(T)$ be the groupoid of line bundles on $T$, with
line-bundle isomorphisms as arrows. A map $f:T'\to T$ gives pullback
$L\mapsto f^*L$. Over $\Spec k$, all line bundles are isomorphic, but the
automorphism group of a one-dimensional vector space is $k^*$. Over $\PP^1$,
the objects include all $\OO_{\PP^1}(d)$. A single isomorphism class over a point
does not make this moduli problem the one-point scheme.

\subsection{Gluing must retain the chosen isomorphisms}
Recall how line bundles glue. For an open cover $\{T_i\}$ of $T$, choose line
bundles $L_i$ and isomorphisms
\[
 \alpha_{ij}:L_j|_{T_i\cap T_j}\xrightarrow{\sim}
                 L_i|_{T_i\cap T_j}.
\]
On triple overlaps require
\[
 \alpha_{ij}\alpha_{jk}=\alpha_{ik}.
\]
Then the $L_i$ glue to a line bundle on $T$, unique up to an isomorphism
compatible with the gluing data. Isomorphisms between two glued bundles also glue.
This is the model for the two stack axioms:
\begin{enumerate}
\item isomorphisms form a sheaf;
\item objects equipped with compatible gluing isomorphisms can be glued.
\end{enumerate}
The second assertion is called \emph{effective descent for objects}. Uniqueness
means uniqueness up to compatible isomorphism; it does not mean that all
automorphism groups vanish.

\begin{definition}[Stack, in working language]
A stack in groupoids assigns to every test scheme $T$ a groupoid $\cX(T)$,
with pullback functors for maps of test schemes, such that isomorphisms and
objects satisfy the two descent conditions above for the chosen topology.
Pullback along a composite agrees with successive pullback through specified
canonical isomorphisms satisfying a compatibility condition.
\end{definition}

The last sentence is the reason textbooks use \emph{categories fibred in
groupoids}. It formalises the fact that $(fg)^*L$ and $g^*f^*L$ are canonically
isomorphic, rather than insisting that they be literally the same set-theoretic
object. No further category theory will be needed for our calculations.

\subsection{Why étale covers enter}
An étale morphism is the algebraic analogue of a map that is locally an
isomorphism without introducing ramification. Formally, it is a flat,
unramified morphism locally of finite presentation. Open immersions are étale,
but there are many other examples.

A model calculation is
\[
 \Spec A[s]/(s^n-a)\longrightarrow\Spec A.
\]
If $a$ is a unit and $n$ is invertible in $A$, then $s$ is a unit upstairs and
the derivative $ns^{n-1}$ is invertible. This is a finite étale cover. It adjoins
an $n$th root of a unit, which generally cannot be chosen on any Zariski open
cover. An étale cover is a jointly surjective family of étale morphisms;
overlaps are fibre products $T_i\times_T T_j$.

In a full definition we can use the fppf topology: covers are jointly
surjective flat morphisms locally of finite presentation. It includes étale
covers. The stacks in these notes satisfy fppf descent. For our smooth groups,
their torsors can be trivialised étale locally, so thinking in terms of étale
covers suffices for the examples. The standard definitions of algebraic and DM
stacks used below are the fppf versions.

\subsection{The simplest twisted family}
Set $T=\Spec k[t,t^{-1}]$ and
\[
 P=\Spec k[s,s^{-1}],\qquad t=s^n.
\]
The group $\mu_n$ acts by $s\mapsto\zeta s$. Every geometric fibre is a free
transitive $\mu_n$-set. This is a \emph{$\mu_n$-torsor}. It is not the trivial
torsor: a section would give an $n$th root of $t$ in $k[t,t^{-1}]$. There is
not even such a root in $k(t)$, as the valuation at $t=0$ would give
$n\operatorname{ord}_0(s)=1$. The same argument excludes a section on any
nonempty Zariski open subset of $T$.

Nevertheless, after the étale base change $P\to T$, the torsor becomes trivial.
This is a concrete reason to enlarge the topology beyond Zariski covers.
It also demonstrates that algebraic closedness of $k$ does not trivialise
torsors over a parameter scheme.

Consider the presheaf sending $T$ to \emph{isomorphism classes} of
$\mu_n$-torsors. The torsor $P$ and the trivial torsor are distinct globally,
but become isomorphic on the cover $P\to T$. Hence that presheaf is not even a
sheaf: distinct sections can become equal on a cover. The groupoid remembers
the gluing isomorphisms that distinguish these families.

\source{For descent and the stack axioms, see Casalaina-Martin--Wise
\cite[\S\S2--3, especially Definition 3.31]{CMW} and Vistoli \cite[Chapters 3--4]{Vistoli}.}

\section{Quotient stacks and classifying stacks}
\label{sec:quotientstacks}

\subsection{A torsor is a family of copies of a group}
A left $G$-torsor $P\to T$ is a scheme with a left $G$-action which is locally
isomorphic, in the chosen topology, to $G\times T$ with translation on $G$.
Equivalently, it is locally nonempty and the map
\[
 G\times P\longrightarrow P\times_T P,\qquad(g,p)\longmapsto(gp,p)
\]
is an isomorphism. In the cases at hand the torsor map is faithfully flat and
locally of finite presentation. A torsor has no preferred identity section;
choosing a section trivialises it.

For $G=\Gm$, torsors are equivalent to line bundles: remove the zero section
from a line bundle, or form the associated line from a torsor. Thus torsors are
already familiar objects with the choice of a local frame left unspecified.

\subsection{Definition of the quotient stack}
\begin{definition}
Suppose $G$ acts on a scheme $X$ on the left. An object of $[X/G](T)$ is a
left $G$-torsor $P\to T$ together with a $G$-equivariant map
\[
 f:P\longrightarrow X,\qquad f(gp)=g f(p).
\]
An arrow is a $G$-equivariant torsor isomorphism commuting with the maps to $X$.
Pullback is obtained by base change.
\end{definition}

One can instead use right torsors and require $f(pg)=g^{-1}f(p)$. Both
conventions give the same stack. We use left torsors here so that all displayed
actions have the same direction.

Why include the torsor? After trivialising it on $T_i$, an object gives maps
$x_i:T_i\to X$. Different trivialisations are related by maps
$g_{ij}:T_i\times_T T_j\to G$, with
\[
 x_i=g_{ij}x_j,\qquad g_{ij}g_{jk}=g_{ik}.
\]
The torsor packages exactly these transition functions. Retaining only global
maps $T\to X$ modulo $G(T)$ would miss the objects with nontrivial transition
functions. The quotient stack is what results from allowing these local
presentations to glue.

Over an algebraically closed field, the torsors for the groups used here are
trivial, so the geometric-point groupoid really is the action groupoid:
\[
 \Aut_{[X/G]}(x)=\Stab_G(x).
\]
This formula is an excellent first calculation, but is not a substitute for
the definition over general $T$.

\begin{example}[The classifying stack]
Take $X=\Spec k$ with trivial $G$-action. Then
\[
 BG=[\Spec k/G],\qquad BG(T)=\{G\text{-torsors over }T\}.
\]
For $G=\mu_n$, there is one geometric isomorphism class, with automorphism
group $\mu_n$. The coarse space will be a point. A family over $T$ still
remembers its torsor, including the nontrivial example of the preceding section.

Equivalently, $B\mu_n(T)$ is the groupoid of pairs
$(N,\phi:N^{\otimes n}\xrightarrow{\sim}\OO_T)$.
The isomorphism $\phi$ is part of the data. Keeping just a line bundle whose
$n$th power happens to be trivial would lose information.
\end{example}

\begin{example}[A free action]
If $X\to Y$ is a $G$-torsor, then $[X/G]\simeq Y$. Indeed, a map $T\to Y$
pulls $X$ back to a torsor over $T$ with an equivariant map to $X$. Conversely,
such a pair descends to a map to $Y$. Thus stacks include ordinary free
quotients. For example,
\[
 [\Gm/\mu_n]\simeq\Gm,\qquad z\longmapsto z^n.
\]
\end{example}

\subsection{Atlases and the first stack fibre product}
There is a natural map $p:X\to[X/G]$, given by the trivial torsor. If
$T\to[X/G]$ corresponds to $(P,f)$, then
\begin{equation}\label{eq:atlasbasechange}
 T\times_{[X/G]}X\simeq P.
\end{equation}
The extra data in this fibre product is a trivialisation of the pulled-back
object, equivalently a point of its torsor. This explains why $p$ is
surjective locally even when an object over $T$ has no global lift to $X$.

If $G$ is smooth, $p$ is a smooth cover. If $G$ is finite étale, $p$ is an
étale cover. These assertions follow directly from
\eqref{eq:atlasbasechange}, since they hold for $P\to T$.

In particular, $\A^1\to[\A^1/\mu_n]$ is étale \emph{at the origin as well}.
Its composite with the coarse map is $z\mapsto z^n$, which is ramified there.
The atlas map and the coarse map have different targets, so there is no
contradiction.

A fibre product of stacks is a \emph{2-fibre product}: an object of
$\cX\times_{\cY}\cZ$ consists of $x$, $z$, and a specified isomorphism between
their images in $\cY$. This is the only extra ingredient beyond the fibre
products of schemes. For the atlas above it gives
\[
 X\times_{[X/G]}X\simeq G\times X.
\]
The two projections are source and target, $x$ and $gx$. The arrows, their
composition, and the object scheme form a \emph{groupoid presentation} of
the stack.
The atlas calculation is a 2-Cartesian square:
\[
\begin{tikzcd}
 P \arrow[r,"f"] \arrow[d] & X \arrow[d,"p"]\\
 T \arrow[r,"{(P,f)}"'] & {[X/G]}.
\end{tikzcd}
\]
The square includes the isomorphism identifying the two induced objects
on $P$. For $T=X$, the square recovers the arrow space $G\times X$.

\source{The torsor construction and its equivalence with the quotient-groupoid
construction are given in the Stacks Project, \href{https://stacks.math.columbia.edu/tag/04UV}{Tag 04UV}.
The algebraicity of smooth groupoid quotients is \href{https://stacks.math.columbia.edu/tag/04TK}{Tag 04TK}.}

\section{The definition of a Deligne--Mumford stack}
\label{sec:dm}

\subsection{A short bridge through algebraic spaces}
An \emph{algebraic space} is a sheaf of sets on the étale site which has a
representable diagonal and an étale surjective cover by a scheme. Here
representability of the diagonal is by schemes. An algebraic space is thus
étale locally a scheme, with ordinary set-valued points rather than groupoid-valued
points. Every scheme gives an algebraic space. The converse is false in general.

This enlargement lets quotients exist in situations where Zariski gluing of
schemes is too restrictive. We do not need an explicit nonschematic algebraic
space in these notes: every displayed coarse space is a scheme. The phrase
``algebraic space'' in the formal definitions prevents an unnecessary
restriction on the general theory.

\subsection{Representability and the diagonal}
A morphism $\cX\to\cY$ is \emph{representable by algebraic spaces} if, after
every map $T\to\cY$ from a scheme, the fibre product $T\times_{\cY}\cX$ is
an algebraic space. A representable morphism is called étale, smooth, finite,
and so on when its base changes to schemes have that property.

The diagonal
\[
 \Delta_{\cX}:\cX\longrightarrow\cX\times\cX
\]
organises isomorphisms. Given two objects $\xi,\eta\in\cX(T)$, its pullback
along $(\xi,\eta):T\to\cX\times\cX$ is
\[
 \underline{\Isom}_T(\xi,\eta).
\]
Over a map $T'\to T$, this sheaf gives isomorphisms between the pulled-back
families. When $\xi=\eta$, it is the automorphism group space of the family.
Thus representability of the diagonal means that isomorphisms between
families themselves form algebraic spaces.

For quotient objects represented by maps $x,x':T\to X$, this is the
transporter
\[
 \{g\in G_T:g x=x'\}.
\]
For the sign action at $x=x'=0$, it is $\mu_2\times T$. For a scheme
considered as a stack, there are only identity arrows: the diagonal expresses
equality, rather than a nontrivial space of isomorphisms.

\begin{definition}[Algebraic stack and DM stack]
An \emph{algebraic stack} is an fppf stack in groupoids whose diagonal is
representable by algebraic spaces and which admits a representable smooth
surjective morphism $U\to\cX$ from a scheme. Such a map is called a smooth atlas.

An algebraic stack is \emph{Deligne--Mumford}, or \emph{DM}, if it admits a
representable étale surjective morphism $U\to\cX$ from a scheme. Such a map
is an étale atlas.
\end{definition}

An étale atlas provides scheme charts with no continuous choices remaining
in each fibre. The stack can still have finite automorphism groups because
different chart identifications are retained as arrows.

We follow the Stacks Project definition. Some books build quasi-compactness
or separatedness conditions on the diagonal into the terminology. Here those
conditions are stated separately. In applications to coarse spaces and descent
below, we impose finite type and separatedness.

\begin{theorem}[Standard criterion; used without proof]\label{thm:dmcriterion}
An algebraic stack is DM if and only if its diagonal is unramified.
In the finite-type quotient situations considered here, over characteristic
zero, $[X/G]$ for a smooth algebraic group $G$ is DM if and only if every
geometric stabiliser is finite.
\end{theorem}

The diagonal formulation says that automorphisms have no infinitesimal motion.
In our finite-type quotient examples, or for the separated finite-type
stacks used below, the stabilisers are finite and reduced. With the broad
Stacks Project convention, DM alone does not impose quasi-compactness
of the diagonal; finiteness requires the additional hypotheses just stated. Over
characteristic zero, finite subgroup schemes of our algebraic groups are
reduced. Outside characteristic zero one must retain the reduced/unramified
condition: a finite group scheme such as $\mu_p$ in characteristic $p$ is
not enough.

\subsection{Examples, nonexamples, and what the definition does not assert}
\begin{enumerate}
\item Every scheme and every algebraic space is a DM stack, with discrete
groupoids of points. An algebraic stack with trivial geometric stabilisers is
an algebraic space; it need not be a scheme.
\item For a finite group $H$ over $k$, the map $X\to[X/H]$ is an étale atlas.
Thus $B\mu_n$, $[\A^1/\mu_n]$, and $[\A^2/\mu_2]$ are DM. The last two
are smooth because they have smooth étale atlases.
\item $B\Gm$ is algebraic, but not DM. Its geometric object has automorphism
group $\Gm$. It is the stack of line bundles introduced earlier.
\item $[\A^1/\Gm]$ for the scaling action is not DM: the origin has stabiliser
$\Gm$, although the nonzero locus has trivial stabiliser.
\item A quotient by a positive-dimensional group can be DM. For example,
$[(\A^2\setminus0)/\Gm]$ with positive weights $(n,1)$ has only finite
stabilisers. Its obvious atlas is smooth of relative dimension one, rather
than étale. Other étale charts exist; we construct them in
Section~\ref{sec:weighted}.
\end{enumerate}

The DM condition does not imply smoothness: take any singular scheme. It
does not imply separatedness: take any nonseparated scheme. Nor does it assert
that a given quotient presentation is by a finite group. It is a local
condition on a stack with an algebraic presentation.

\paragraph{Separatedness.}
A stack is separated over $k$ if its diagonal is proper. For a finite-type
DM stack the diagonal is already unramified and quasi-finite; if it is proper
it is finite. This explains why finite stabilisers and well-behaved finite
quotient charts occur together in separated DM geometry. Closed orbits and
separation questions must still be checked in a GIT application; finite
stabilisers alone are not a general substitute for those checks.

\paragraph{Dimension and smoothness.}
For a smooth group $G$ acting on $X$, the smooth atlas has relative dimension
$\dim G$, so $\dim[X/G]=\dim X-\dim G$ in the equidimensional examples.
The quotient stack is smooth over $k$ when $X$ is smooth: smoothness can be
checked on a smooth atlas. Thus $\PP(n,1)$ is a smooth curve, while
$[\A^2/\mu_2]$ is a smooth surface.

\paragraph{Relation to orbifolds.}
Over $\mathbb C$, a smooth separated DM stack behaves locally like a complex
manifold modulo a finite group. This is the useful orbifold picture. Some
definitions of orbifold require every local group to act faithfully. DM
stacks allow ineffective actions: $B\mu_n$ is a simple example, since its
tangent space is zero. A singular DM stack is also outside the smooth
orbifold picture. We will therefore use the algebraic definition rather than
identify the two terms without qualifications.

\source{Definitions: Stacks Project, \href{https://stacks.math.columbia.edu/tag/026N}{Tag 026N}.
The diagonal criterion is \href{https://stacks.math.columbia.edu/tag/06N3}{Tag 06N3}.
For the diagonal as an isomorphism sheaf, see \cite[Example 4.6]{CMW}.}

\section{Coarse spaces and local finite quotients}
\label{sec:coarse}

\begin{definition}[Coarse moduli space]
A morphism $\pi:\cX\to X$ to an algebraic space is a coarse moduli space if:
\begin{enumerate}
\item for every algebraically closed extension $\Omega/k$, it induces a
bijection between isomorphism classes in $\cX(\Omega)$ and $X(\Omega)$;
\item it is universal for morphisms from $\cX$ to algebraic spaces: every
such morphism factors uniquely through $X$.
\end{enumerate}
\end{definition}

Both clauses matter. The first says that geometric isomorphism classes are
preserved as points. The second controls morphisms and functions. Neither
says that families over general schemes are preserved.

For a finite group $H$ acting on an affine scheme $\Spec A$, the coarse space
of $[\Spec A/H]$ is $\Spec A^H$. A morphism from the stack to a scheme is
the same as an $H$-invariant morphism from the atlas, so the invariant-ring
calculation is consistent with the universal property. The universality for
general algebraic-space targets is part of the quotient theorem.

\begin{theorem}[Coarse spaces and local structure; used without proof]
A separated finite-type DM stack over $k$ admits a coarse moduli space,
which is an algebraic space. Étale locally on this coarse space it has
finite-group quotient presentations. At a chosen geometric point, a chart
can be chosen with the stabiliser of that point as the local finite group.
\end{theorem}

This gives a way to learn the general theory through finite quotient
calculations. It is not a statement that every stack has one global
presentation by a finite group, nor that its coarse space must be a scheme.

\subsection{A coarse map preserves points but changes families}
The coarse space of $B\mu_n$ is $\Spec k$. There is one geometric
isomorphism class on each side. But a map $T\to B\mu_n$ carries a torsor,
whereas a map $T\to\Spec k$ carries no such data. The nontrivial torsor
$s^n=t$ and the trivial torsor are sent to the same map to the coarse point.

The coarse space of $[\A^1/\mu_n]$ is $\A^1_u$. Its nontrivial stabiliser
is invisible in that smooth line. In the quadratic-cone example, the coarse
space instead has a singularity. Thus the implication
``smooth stack $\Rightarrow$ smooth coarse space'' is false, and
``smooth coarse space $\Rightarrow$ no stack structure'' is also false.

\subsection{Good moduli spaces are a different notion}
For a reductive GIT quotient with strictly semistable points, different
orbits can have closures meeting the same closed orbit. The GIT quotient
identifies them. Such a quotient need not be a coarse moduli space of the
whole semistable quotient stack.

The smallest example is $[\A^1/\Gm]$ for scaling. It has two geometric
isomorphism classes: the zero orbit and the nonzero orbit. Its affine
invariant quotient is $\Spec k$, since $k[z]^{\Gm}=k$. That point is
\emph{not} its coarse moduli space: it fails the bijection in the definition.

This quotient is a \emph{good moduli space}. For our purposes its defining
properties are that the pushforward on quasi-coherent sheaves is exact and
$\OO_X\simeq\pi_*\OO_{\cX}$; the morphism is also required to be
quasi-compact. These conditions extend the behaviour of invariant rings and
invariant modules to stacks. We will use the term only to distinguish the
quotients at critical cutting levels from the coarse quotients at regular ones.

\source{Coarse-space existence is discussed in the Stacks Project,
\href{https://stacks.math.columbia.edu/tag/04UX}{Tag 04UX}; local structure is
discussed in \cite[\S3.4]{Behrend}. For good moduli spaces see
\cite[Definition 4.1 and \S13]{Alper}.}

\section{Sheaves and line bundles retain stabiliser representations}
\label{sec:sheaves}

\subsection{Equivariant sheaves are sheaves on the quotient stack}
A vector bundle on $[X/G]$ is equivalent to a $G$-equivariant vector bundle
on $X$. An equivariant structure says how $g$ identifies the fibre at $x$
with the fibre at $gx$, compatibly with multiplication in $G$. This is the
descent datum on the atlas and its arrow space $G\times X$.

At a point fixed by $g$, the equivariant structure is an action on the same
fibre. Hence a vector bundle on a stack gives a representation of every
stabiliser. In rank one this representation is a \emph{character}.
These characters are additional geometric data, even if the underlying
bundle on the atlas is trivial.

\begin{example}[All line bundles on $B\mu_n$]
A vector bundle on $B\mu_n$ is a representation of $\mu_n$. Its line bundles
are the one-dimensional representations
\[
 \chi^c:\mu_n\to\Gm,\qquad\zeta\longmapsto\zeta^c,
 \qquad c\in\mathbb Z/n\mathbb Z.
\]
Thus $\Pic(B\mu_n)\simeq\mathbb Z/n\mathbb Z$, while the coarse point has
trivial Picard group. The global sections of the character line are its
invariant vectors: $k$ for $c=0$, and zero for $c\ne0$.
\end{example}

\subsection{Character conventions and invariant sections}
On $[\A^1_z/\mu_n]$, let $\cL_c$ be the trivial line upstairs with the
\emph{geometric fibre action}
\[
 (z,v)\longmapsto(\zeta z,\zeta^c v).
\]
A section descends to the stack if it is equivariant. Writing it as the
function $f(z)$ in this trivialisation gives
\begin{equation}\label{eq:equivsection}
 f(\zeta z)=\zeta^c f(z).
\end{equation}
Thus the allowable monomials are $z^j$ with $j\equiv c\pmod n$.

This equation avoids an ambiguity about representations on functions. With
the inverse-pullback convention $g\cdot f=f\circ g^{-1}$, the coordinate $z$
has weight $-1$. With the direct-pullback convention $f\mapsto f\circ g$ for
our abelian groups, it has weight $+1$. The \emph{geometric tangent action}
$z\mapsto\zeta z$ is fixed throughout. Always state which action defines a
weight before comparing formulas.

\subsection{Pushforward is not descent}
Let $\pi:[\A^1/\mu_n]\to\A^1_u$, $u=z^n$. Choose $r$ with
$0\leq r<n$ and $r\equiv c\pmod n$. Equation~\eqref{eq:equivsection} gives
\begin{equation}\label{eq:rankonepush}
 \Gamma(\A^1_u,\pi_*\cL_c)=z^r k[u].
\end{equation}
As a $k[u]$-module this is free of rank one. Therefore the pushforward is an
ordinary line bundle, even when $c\not\equiv0\pmod n$.

Nevertheless $\cL_c$ does not descend from the coarse space when
$c\not\equiv0\pmod n$.
If $\cL_c\simeq\pi^*L$, the stabiliser at $z=0$ would act trivially on its
fibre: the coarse point has no nontrivial automorphisms that could act on $L$.
But that fibre has character $c$. The natural comparison
\[
 \pi^*\pi_*\cL_c\longrightarrow\cL_c
\]
is locally given by multiplication by $z^r$ in the displayed trivialisation.
For $r>0$ it vanishes at $0$, so it is not an isomorphism.

There are therefore three different operations or questions:
\begin{enumerate}
\item the original line bundle $\cL$ on the stack;
\item its sheaf $\pi_*\cL$ on the coarse space;
\item whether $\cL$ is isomorphic to the pullback of a line bundle there.
\end{enumerate}
Even an invertible pushforward does not answer the third question.

\begin{theorem}[Descent criterion in our setting; used without proof]
Let $\pi:\cX\to X$ be the coarse map of a separated finite-type DM stack
over $k$. A vector bundle $E$ on $\cX$ is pulled back from a vector bundle
on $X$ if and only if every geometric stabiliser acts trivially on its fibre.
In particular this applies to line bundles.
\end{theorem}

Characteristic zero makes the finite stabilisers linearly reductive, so these
coarse maps are tame good moduli spaces. The necessity of the criterion was
proved above. The sufficiency is a general descent theorem, rather than a
consequence of counting sections. In a finite quotient presentation it says
that equivariant gluing with trivial fibre actions is ordinary quotient
gluing. See Alper's vector-bundle descent theorem \cite[Theorem 10.3]{Alper}.

\subsection{An invariant module which is not a line bundle}
Return to $\cX=[\A^2_{x,y}/\mu_2]$ with simultaneous sign action, and give a
trivial line upstairs the sign fibre character. Set
\[
 A=k[x^2,xy,y^2],\qquad
 M=\{f\in k[x,y]:f(-x,-y)=-f(x,y)\}.
\]
The pushforward is the coherent sheaf associated to the $A$-module $M$.
It is generated by $x,y$. At the cone vertex $\mathfrak m=(x^2,xy,y^2)$,
the vector space $M/\mathfrak m M$ has dimension two: the classes of $x,y$
remain independent. On the other hand $M$ has rank one over the fraction
field of $A$, because $y=(xy/x^2)x$ generically. A free rank-one module at
the vertex would have a one-dimensional quotient modulo $\mathfrak m$.
Thus $M$ is not locally free there.

This is a concrete model for why the pushforward of an honest stack line
bundle can fail to be invertible on the coarse space. It complements
\eqref{eq:rankonepush}, where the pushforward is invertible but descent
still fails.

\paragraph{Why invariants behave well here.}
For a finite group $H$ in characteristic zero, averaging
$v\mapsto |H|^{-1}\sum_{h\in H}h v$ projects onto invariant vectors.
Consequently taking invariants is exact. On affine finite-quotient charts,
pushforward is the corresponding invariant-module operation. Exactness is
very useful, but it does not make pushforward preserve tensor products, as
the next example will show.

\source{Equivariant sheaves and invariant sections are explained in
\cite[\S3.5, especially Exercise 3.56]{Behrend}.}

\section{A weighted projective line, worked out in charts}
\label{sec:weighted}

Let $\Gm$ act on $U=\A^2_{x,y}\setminus\{0\}$ by
\begin{equation}\label{eq:weightedaction}
 t\cdot(x,y)=(t^n x,t y).
\end{equation}
Define
\[
 \cP=\PP(n,1):=[U/\Gm].
\]
This is the same weighted projective line often written $\PP(1,n)$ after
interchanging the coordinates. We retain the order $(x,y)$ to match the cut
in Section~\ref{sec:toycut}.

If $y\ne0$, a stabilising $t$ satisfies $t=1$. At $(x,0)$ with $x\ne0$,
it satisfies $t^n=1$. All stabilisers are finite, so $\cP$ is DM. It is
smooth because $U$ is smooth.

\subsection{The two charts}
On $y\ne0$, there is a unique way to set $y=1$: choose $t=y^{-1}$. The
remaining coordinate is
\[
 z=\frac{x}{y^n}.
\]
Thus this open substack is an ordinary affine line $\A^1_z$.

On $x\ne0$, setting $x=1$ requires a choice $t^n=x^{-1}$. Such a choice
exists étale locally, but has a $\mu_n$ ambiguity. The remaining coordinate
$s$ transforms by $s\mapsto\zeta s$. Consequently this chart is
\[
 [\A^1_s/\mu_n].
\]
Its coarse coordinate is
\[
 u=s^n=\frac{y^n}{x}.
\]
On the overlap $zu=1$. The two coarse affine lines therefore glue to
$\PP^1$. The coarse map is
\begin{equation}\label{eq:coarseweighted}
 \pi:\cP\longrightarrow\PP^1,\qquad [x:y]\longmapsto[x:y^n].
\end{equation}
The expression on the right is unchanged as a projective point when
$(x,y)$ is replaced by $(t^n x,ty)$.

Taking the ordinary affine line atlas of the second quotient chart and the
first affine chart gives an étale atlas of $\cP$. This constructs the atlas
whose existence was asserted by the DM criterion. It is different from
the smooth $\Gm$-torsor atlas $U\to\cP$.

\subsection{A description in families}
An object of $\cP(T)$ can be described as
\[
 (N,\,x\in\Gamma(T,N^{\otimes n}),\,
        y\in\Gamma(T,N)),
\]
where the two sections do not vanish simultaneously at any geometric point.
An isomorphism of line bundles must carry the displayed sections to the
corresponding sections. This is the family version of weighted coordinates.
In the left-torsor convention above, the corresponding torsor is
$P=(N^\vee)^\times$: a nonzero covector evaluates the sections to give
the coordinates. Scaling that covector by $t$ multiplies those
coordinates by $t^n$ and $t$, as required.

For example, where $y$ is nonzero it trivialises $N$, leaving the ordinary
function $x/y^n$. Where $x$ is nonzero it trivialises $N^{\otimes n}$,
leaving an $n$th-root choice for $N$. Thus the chart calculation reflects
actual families, rather than merely orbit notation.

\subsection{The boundary and the line bundles}
Let $\cD=\{y=0\}\subset\cP$. It is
\[
 \cD=[\Gm/\Gm]\simeq B\mu_n,
\]
where the action on the first $\Gm$ is $t\cdot x=t^n x$. Its coarse space
is the point $D=[1:0]\in\PP^1$.

Define $\OO_{\cP}(c)$ using the geometric fibre action
\[
 (x,y,v)\longmapsto(t^n x,ty,t^c v).
\]
Its global sections are weighted homogeneous polynomials of degree $c$.
The coordinate $y$ is an equivariant section of $\OO_{\cP}(1)$ and vanishes
simply along $\cD$. Hence
\begin{equation}\label{eq:weightlines}
 \OO_{\cP}(1)\simeq\OO_{\cP}(\cD),\qquad
 \pi^*\OO_{\PP^1}(1)\simeq\OO_{\cP}(n).
\end{equation}
The second identity can also be read directly from the two degree-$n$
sections $x,y^n$ defining the coarse map.

At $\cD$, the stabiliser acts on $\OO_{\cP}(c)$ by $\zeta^c$.
It follows that
\begin{equation}\label{eq:weighteddescent}
 \OO_{\cP}(c)\text{ descends to }\PP^1
 \quad\Longleftrightarrow\quad n\mid c.
\end{equation}
Sufficiency follows already from \eqref{eq:weightlines}; necessity follows
from the fibre character. We have therefore verified this case of the
descent theorem without using its general proof.

\subsection{Counting sections and computing pushforward}
For $m\ge0$, the monomials of degree $m$ are
\[
 x^a y^{m-na},\qquad 0\leq a\leq\lfloor m/n\rfloor.
\]
Consequently
\begin{equation}\label{eq:weightedh0}
 \dim H^0(\cP,\OO_{\cP}(m))=\lfloor m/n\rfloor+1.
\end{equation}
No higher-cohomology machinery is needed: this just counts equivariant
polynomials on $U$. A regular function on $\A^2\setminus0$ extends to
$\A^2$; explicitly, $k[x,y,y^{-1}]\cap k[x,y,x^{-1}]=k[x,y]$ in $k(x,y)$.

More precisely, for every $m\in\mathbb Z$,
\begin{equation}\label{eq:weightedpush}
 \pi_*\OO_{\cP}(m\cD)
 \simeq\OO_{\PP^1}\bigl(\lfloor m/n\rfloor D\bigr).
\end{equation}
To see this locally near $D$, use a root coordinate $s$ with $u=s^n$.
The divisor sheaf upstairs is $s^{-m}k[s]$, and its invariant rational
monomials are $s^{nj}=u^j$ with $nj\ge-m$. Thus
\[
 (s^{-m}k[s])^{\mu_n}=u^{-\lfloor m/n\rfloor}k[u].
\]
Off $D$ both divisor sheaves are canonically trivial. This proves the formula
by gluing and includes negative $m$.

When $0<m<n$, the pushforward in \eqref{eq:weightedpush} is $\OO_{\PP^1}$,
but the stack line does not descend. In particular,
\[
 (\pi_*\OO_{\cP}(1))^{\otimes n}\simeq\OO_{\PP^1},\qquad
 \pi_*(\OO_{\cP}(1)^{\otimes n})\simeq\OO_{\PP^1}(1).
\]
Pushforward has forgotten stabiliser characters, and it does not commute
with this tensor product.

\reading{For $n=3$, write down the sections of $\OO(m)$ for $0\le m\le6$.
Which powers descend? Can the dimensions of sections alone recover the
character at the boundary?}

\section{Root stacks, divisors, and normal lines}
\label{sec:roots}

\subsection{Making the root of a divisor into geometric data}
Let $D\subset X$ be an effective Cartier divisor. Its line bundle $\OO_X(D)$
has a canonical section $s_D$ with zero divisor $D$. An $n$th root of this
pair should be a line bundle $N$ with a section $s$ and an isomorphism
$N^{\otimes n}\simeq\OO_X(D)$ carrying $s^n$ to $s_D$. Such a pair
need not exist on $X$. Rather than choose one, we parameterise all choices.

\begin{definition}[Divisor root stack]
An object of $\sqrt[n]{(X,D)}$ over a scheme $T$ is a map $f:T\to X$ and
data
\[
 (N,s,\phi),\qquad
 \phi:N^{\otimes n}\xrightarrow{\sim}f^*\OO_X(D),\qquad
 \phi(s^{\otimes n})=f^*s_D.
\]
Arrows are line-bundle isomorphisms compatible with $s$ and $\phi$.
\end{definition}

If $D$ has equation $a\in A$ on $\Spec A$, its local presentation is
\begin{equation}\label{eq:rootchart}
 \sqrt[n]{(\Spec A,D)}
 \simeq\left[\Spec A[z]/(z^n-a)\,/\mu_n\right],\qquad
 \zeta\cdot z=\zeta z.
\end{equation}
Changing a trivialisation changes the root coordinate and glues these
quotient charts. Since $n$ is invertible, this is a DM stack. If $X$ and
$D$ are smooth, it is smooth: locally $a$ is a regular coordinate and the
equation $a=z^n$ replaces it by the smooth coordinate $z$. For singular
$D$, smoothness is not automatic.

Away from $D$, the section $s$ is nonzero and trivialises $N$, so the root
stack is just $X\setminus D$. Over $D$, scaling the root by $\zeta$ preserves
its $n$th power. The stabiliser is $\mu_n$. Locally the invariant ring in
\eqref{eq:rootchart} is $A$, so the coarse space is $X$.

Our weighted line is the example
\[
 \PP(n,1)\simeq\sqrt[n]{(\PP^1,D)},\qquad D=[1:0].
\]
Its two charts are exactly the ordinary chart away from $D$ and
$[\A^1_s/\mu_n]$ near $D$.

\subsection{The reduced root divisor and the scheme-theoretic fibre}
The universal root section vanishes along a Cartier substack $\cD$ with
\[
 \pi^*D=n\cD,\qquad
 \OO(n\cD)\simeq\pi^*\OO_X(D).
\]
If $D$ is smooth, this $\cD$ is the reduced substack over $D$. Its local
description is $[D/\mu_n]$ with trivial action on $D$. Globally it can be
a nontrivial $\mu_n$-gerbe: an object over $T\to D$ is an $n$th root of
$\OO_X(D)|_T$. This is locally $D\times B\mu_n$, but a global product
would require a compatible global root. When $D$ is the single $k$-point
in our weighted-line example, $\cD\simeq B\mu_n$.

There is a useful scheme-theoretic subtlety. For $[\A^1_z/\mu_n]\to\A^1_u$
with $u=z^n$, the actual fibre over $u=0$ is
\[
 \left[\Spec k[z]/(z^n)\,/\mu_n\right].
\]
It is a nonreduced thickening of $B\mu_n$, rather than just $B\mu_n$.
The reduced residual gerbe, cut out by $z=0$, is $B\mu_n$.
The distinction is invisible on algebraically closed field-valued points
but visible on nonreduced test schemes. It is another reason to retain
the full functor of points.

\subsection{Normal and conormal characters}
At the origin of the chart $\A^1_z$, a normal vector transforms as
$v\mapsto\zeta v$. The normal line of $B\mu_n$ in $[\A^1/\mu_n]$
therefore has character $+1$. The conormal line is its dual and has
character $-1$. This statement concerns the geometric action on vectors
and covectors, independently of which convention is used to write an
action on functions.

If $i:\cD\hookrightarrow\cX$ is a Cartier divisor, its normal line is
\[
 \normal_{\cD/\cX}\simeq\OO_{\cX}(\cD)|_{\cD}.
\]
For the weighted line this is the character-$1$ line on $B\mu_n$.
The coarse divisor $D\subset\PP^1$ also has an ordinary normal line.
Its pullback to $B\mu_n$ has trivial stabiliser character, so it is not
the stack normal line. Instead,
\begin{equation}\label{eq:normalroot}
 \normal_{\cD/\cX}^{\otimes n}
 \simeq (\pi|_{\cD})^*\normal_{D/X}.
\end{equation}
This follows by restricting $\OO(n\cD)=\pi^*\OO(D)$.
Thus even when both spaces and their coarse divisors are smooth, the
normal line on the stack can carry information lost by the coarse pair.

The line $\OO(\cD)$ is sometimes thought of as the line of the
``fractional divisor'' $D/n$. It is an honest line bundle on the stack.
The fractional notation on the coarse space is a way of remembering the
root, not an assertion that $D/n$ is an integral Cartier divisor there.

\paragraph{Root stacks versus root gerbes.}
If one asks only for $N^{\otimes n}\simeq L$ and does not include a
section, the resulting root \emph{gerbe} has $\mu_n$ stabiliser everywhere.
The divisor root stack includes a section and has trivial stabiliser off
its zero locus. For instance, the root gerbe of the trivial line on $X$
is $X\times B\mu_n$, whereas $\sqrt[n]{(\A^1,0)}$ has stabiliser only
at the origin. Here $0$ in the latter notation means the origin divisor.

\source{The divisor root construction is due to Cadman \cite[\S2]{Cadman}.
For the concrete curve description and divisor pushforward formula, see
\cite[\S3.6, Definition 3.62 and Theorem 3.64]{Behrend}.}

\section{Where these ideas enter algebraic cutting}
\label{sec:cutting}

\subsection{The setting and the chosen sign convention}
Suppose $M$ is a smooth variety with a $\Gm$-action and an ample
$\Gm$-linearised line bundle $L$. A linearisation specifies maps on fibres
$t_L:L_p\to L_{tp}$ compatible with the group action. The section ring is
\[
 R=\bigoplus_{m\ge0}H^0(M,L^{\otimes m}),\qquad
 R_{m,w}=H^0(M,L^{\otimes m})_w.
\]
The index $m$ is a tensor or projective degree, while $w$ is a group weight.
They are independent. For semiprojective examples the degree-zero ring
$R_0=H^0(M,\OO_M)$ can be nonconstant.

For this discussion assume $R$ is finitely generated, $M\simeq\Proj R$,
and the affinisation $M\to\Spec R_0$ is projective. A contracting action
is an additional condition, not a consequence of being projective over an
affine base. Our toy model $M=\A^1$ has both properties.

We use the convention in the repository's
\path{claude-notes/conventions.tex}, dated 12 November 2026. For sections,
\[
 (t\cdot s)(p)=t_L^{-1}s(tp).
\]
For functions this is direct pullback, $f(p)\mapsto f(tp)$. Because our
group is abelian, these operators form a representation. A fixed component
$F$ has section weight $\Phi(F)$, so the \emph{geometric fibre weight}
of $L|_F$ is $-\Phi(F)$.

Introduce an auxiliary coordinate $y$ with point action
$y\mapsto t^\epsilon y$, where $\epsilon\in\{+1,-1\}$. Twist by a trivial
line with geometric fibre weight $d\in\mathbb Z$, whose constant section
$\one_d$ has section weight $-d$. On
$X_\epsilon=M\times\A^1_y$, the section
\[
 r\,y^\ell\one_d^m,\qquad r\in R_{m,w},
\]
is invariant exactly when
\begin{equation}\label{eq:cutinvariance}
 w+\epsilon\ell-md=0,\qquad \ell\ge0.
\end{equation}
It follows that the invariant section algebra is
\begin{align}
 B^+_m&=\bigoplus_{w\le md}R_{m,w}y^{md-w}\one_d^m,\label{eq:cutplus}\\
 B^-_m&=\bigoplus_{w\ge md}R_{m,w}y^{w-md}\one_d^m.\label{eq:cutminus}
\end{align}
These equations, rather than an unstated moment-map sign, determine the
two directions in this note.

\subsection{Stack cut, coarse cut, and boundary}
Let $X^{\sst}_\epsilon(d)$ be the locus on which some invariant section
of positive tensor degree is nonzero. The \emph{stack cut} and its GIT
quotient are
\[
 \cM^{\leq d}=[X^{\sst}_+(d)/\Gm],\quad
 M^{\leq d}=\Proj B^+,
\]
and similarly $\cM^{\geq d}$ and $M^{\geq d}$ using $\epsilon=-1$.
At regular levels in the setting below, the GIT quotients are coarse
spaces. At a general level they are good moduli spaces.

Putting $y=0$ in \eqref{eq:cutinvariance} leaves $w=md$. The common
boundary is therefore described by
\[
 \cM_d=[M^{\sst}(d)/\Gm],\qquad
 M_d=\Proj\bigoplus_mR_{m,md}.
\]
The stack cut carries the line bundle $\cL$ defined by the product
linearisation. The corresponding graded sheaf $\widetilde{B^\epsilon(m)}$
on the quotient is its pushforward in this setting. It need not be
invertible, and invertibility alone would still not prove that the stack
line descends. Section~\ref{sec:sheaves} gives both kinds of example.

\subsection{Why the stabilisers are concentrated on the boundary}
A point $(p,y)$ is fixed by $t$ only if
\[
 tp=p,\qquad t^\epsilon y=y.
\]
If $y\ne0$ and $\epsilon=\pm1$, the second equation forces $t=1$.
If $y=0$, the stabiliser is $\Stab_{\Gm}(p)$. Thus
\begin{equation}\label{eq:cutstabs}
 \Stab(p,y)=
 \begin{cases}
 \{1\},&y\ne0,\\
 \Stab_{\Gm}(p),&y=0.
 \end{cases}
\end{equation}
Closed algebraic subgroups of $\Gm$ over $k$ are $\Gm$ itself or finite
cyclic groups $\mu_r$ (including the trivial group).

In the current repository setting, a level is called regular when
$d\ne\Phi(F)$ for all fixed components. Then the semistable locus has
finite stabilisers and closed orbits within that locus. Its quotient stack
is smooth DM and has a geometric coarse quotient. Regularity permits
finite cyclic stabilisers on the boundary; it does not mean that the action
is free. At a critical level, fixed points can remain semistable and yield
$\Gm$ stabilisers. The resulting quotient stack is then algebraic but
not DM along those points.

One can already see the critical equality from section weights. At a fixed
point $p\in F$, a nonvanishing section of $L^{\otimes m}$ has weight
$m\Phi(F)$. At $y=0$, invariance demands $m\Phi(F)=md$. Thus fixed
points can survive on the boundary exactly at the critical weight.

\subsection{Normal character versus polarisation character}
On the regular boundary the auxiliary normal direction transforms by
$t^\epsilon$. The stack normal line is the associated line
\begin{equation}\label{eq:cutnormal}
 \normal_{\cM_d/\cM^{\leq_\epsilon d}}
 =[(M^{\sst}(d)\times\charline{\epsilon})/\Gm].
\end{equation}
Here $\leq_\epsilon$ means $\leq$ for $\epsilon=+1$ and $\geq$ for
$\epsilon=-1$. Formula~\eqref{eq:cutnormal} follows by taking the normal
bundle of $y=0$ before quotienting. It does not require a global
identification of a whole boundary neighbourhood with a normal bundle.

If the boundary stabiliser is $\mu_r$, the two normal characters are
$\zeta$ and $\zeta^{-1}$. They are dual lines. The polarisation is instead
determined by $L$ and the level twist. There is no general reason for it
to equal the normal line. The ordinary coarse normal sheaf is a third
object; the root example already proves that pulling it back need not
give the stack normal line.

\paragraph{A distinction from residual symmetry.}
The stabiliser in \eqref{eq:cutstabs} belongs to the group being
\emph{quotiented}. A separate residual torus can act on the resulting cut.
A fixed point of that residual action is not automatically an additional
automorphism of the object in the stack. Isotropy of the quotient stack
and fixed loci of a residual action are different data.

\section{One complete cutting toy model}
\label{sec:toycut}

Let
\[
 M=\A^1_x,\qquad t\cdot x=t^n x,\qquad L=\OO_M
\]
with trivial geometric fibre action on $L$. This is an affine example:
its affinisation is the identity morphism and is therefore projective.
The trivial bundle is ample on
the affine scheme. The action contracts to $0$. The only fixed component
is $0$, with $\Phi(0)=0$; the level $d=1$ is regular.

If $h$ denotes the constant degree-one section of $L$, then
$R=k[x,h]$ with tensor degrees $\deg x=0,\deg h=1$ and section weights
$\operatorname{wt}x=n,\operatorname{wt}h=0$. In particular, this is an
example where the degree-zero ring is nonconstant. Suppressing $h$ and
retaining the twist symbol, invariant sections of tensor degree $m$ have
the form $x^a y^b\one_1^m$ with
\begin{equation}\label{eq:toyinvariance}
 na+\epsilon b=m,\qquad a,b\ge0.
\end{equation}

\subsection{The lower cut is the weighted projective line}
For $\epsilon=+1$, the invariant sections are nonzero somewhere whenever
$(x,y)\ne(0,0)$: if $y\ne0$ use $y\one_1$; if $x\ne0$ use
$x\one_1^n$. All positive-degree sections vanish at the origin. Hence
\[
 X^{\sst}_+(1)=\A^2\setminus0,
 \qquad\cM^{\le1}=\PP(n,1).
\]
The invariant algebra is a polynomial ring in
\[
 A=x\one_1^n,\qquad B=y\one_1,
 \qquad \deg_m A=n,\quad\deg_m B=1.
\]
Its $n$th Veronese is $k[A,B^n]$; after regrading, both generators have
degree one. Therefore
\[
 M^{\le1}=\Proj k[A,B]\simeq\PP^1.
\]
This computes the stack cut, coarse cut, and grading in the same example.

Where $y\ne0$ the quotient has coordinate $x/y^n$ and is $\A^1$.
The added boundary is $B\mu_n$, with coarse point $[1:0]$. Its normal
character is $+1$. The cut polarisation has fibre character $+1$ as well
in this particular example, so it is $\OO_{\PP(n,1)}(1)$ and does not
descend when $n>1$. Its $n$th power does descend.

\subsection{The upper cut is an affine root stack}
For $\epsilon=-1$, equation~\eqref{eq:toyinvariance} is $na-b=m$.
Positive tensor degree forces $a>0$, so $x=0$ is unstable. Conversely,
if $x\ne0$, the section $x\one_1^n$ is nonzero. Thus
\[
 X^{\sst}_-(1)=\mathbb G_{m,x}\times\A^1_y,
 \qquad
 \cM^{\ge1}=[(\Gm\times\A^1)/\Gm],
\]
where $t\cdot(x,y)=(t^n x,t^{-1}y)$. Setting $x=1$ étale locally
leaves $\mu_n$ acting on the normal coordinate by $y\mapsto\zeta^{-1}y$:
\[
 \cM^{\ge1}\simeq[\A^1_y/\mu_n],\qquad
 M^{\ge1}\simeq\A^1_v,\quad v=xy^n.
\]
For a ring check, the $n$th Veronese of $B^-$ is generated over $k[v]$
by $A=x\one_1^n$ in regraded degree one. Indeed a monomial in degree
$nr$ is $x^{r+j}y^{nj}\one_1^{nr}=A^r v^j$.
Hence its projective spectrum is $\Spec k[v]$.

The interior $y\ne0$ is $\mathbb G_{m,v}$, while the boundary is again $B\mu_n$
at $v=0$. This upper cut is not proper over $k$. Its stack normal line
has character $-1$, the dual of the lower-cut normal line. The same
polarisation twist has character $+1$ on both boundary stacks. Thus on
the upper cut it is the dual of the boundary normal line.

\begin{center}
\begin{tabular}{@{}lll@{}}
\toprule
 & Lower cut & Upper cut\\
\midrule
Semistable locus & $\A^2\setminus0$ & $\Gm\times\A^1$\\
Stack & $\PP(n,1)$ & $[\A^1/\mu_n]$\\
Coarse space & $\PP^1$ & $\A^1$\\
Interior & $\A^1$ & $\Gm$\\
Boundary stack & $B\mu_n$ & $B\mu_n$\\
Normal character & $+1$ & $-1$\\
Boundary polarisation character & $+1$ & $+1$\\
\bottomrule
\end{tabular}
\end{center}
The boundary identifications use the same subgroup $\mu_n\subset\Gm$.
Keeping this identification is what makes the two normal characters
meaningfully opposite.

\subsection{Weights, sections, and the boundary contribution}
On $M$, the monomial $x^a$ has weight $na$. The lower-cut degree-$m$
sections are those with $na\le m$:
\[
 H^0(\cM^{\le1},\cL^{\otimes m})
 \simeq\operatorname{span}_k\{1,x,\ldots,x^{\lfloor m/n\rfloor}\}
 \quad(m\ge0).
\]
Explicitly $x^a$ on the right corresponds to
$x^a y^{m-na}\one_1^m$ on the left. The displayed isomorphism is the
polynomial identification of these vector spaces; it keeps the original
residual weights $na$ if the residual action scales $x$ alone.

On the boundary, $\cL^{\otimes m}$ has character $m\bmod n$, so
\begin{equation}\label{eq:boundarysections}
 H^0(B\mu_n,\cL^{\otimes m}|_{B\mu_n})
 =\begin{cases}k,&n\mid m,\\0,&n\nmid m.\end{cases}
\end{equation}
For $n=2$, this gives the table
\[
\begin{array}{c|rrrrrrr}
 m&0&1&2&3&4&5&6\\\hline
 \dim H^0(\cM^{\le1},\cL^m)&1&1&2&2&3&3&4\\
 \dim H^0(B\mu_2,\cL^m)&1&0&1&0&1&0&1.
\end{array}
\]
For $m\ge1$, multiplication by $y$ embeds the degree-$(m-1)$
polynomials into the degree-$m$ polynomials. The quotient has a monomial
with $y$-exponent zero exactly when $n\mid m$, namely $x^{m/n}$.
This proves the boundary contribution directly, without invoking any
cohomological vanishing theorem. On sheaves the underlying divisor
sequence is
\[
 0\to\OO_{\cP}(m-1)\xrightarrow{\ y\ }
 \OO_{\cP}(m)\to\OO_{\cD}(m)\to0.
\]

Replacing the boundary stack by its coarse point and replacing its
character line by an ordinary one-dimensional fibre would incorrectly
give one boundary section for \emph{every} $m$. The periodic behaviour in
\eqref{eq:boundarysections} is a first elementary trace of the
representation data relevant to orbifold index calculations.

\subsection{What changes at the critical level}
Keep the same $M,L$, take $d=0$, and choose $\epsilon=+1$. Invariance is
$na+b=0$, so the only invariant sections are constant in $x,y$.
Their positive tensor powers are nonzero everywhere. Therefore
\[
 X^{\sst}_+(0)=\A^2,\qquad
 \cM^{\le0}=[\A^2/\Gm],\qquad M^{\le0}=\Spec k.
\]
The origin has stabiliser $\Gm$, so this stack is not DM. It also has
many distinct geometric orbits mapping to the same point. Thus the point
is a good moduli space, not a coarse moduli space of the whole quotient
stack. This calculation isolates the two changes at a wall: continuous
stabilisers can appear, and the GIT quotient can cease to classify
geometric isomorphism classes.

\section{Exercises with solutions}
\label{sec:exercises}

\begin{exercise}\label{ex:trivial}
Compare $[\A^1/\mu_n]$ for the trivial action and for scalar multiplication.
Find their coarse spaces, stabilisers, and the open substacks away from zero.
\end{exercise}
\begin{exercise}\label{ex:torsor}
Show that $[\Gm/\mu_n]\simeq\Gm$ for scalar multiplication, and explain why
this does not extend to an isomorphism $[\A^1/\mu_n]\simeq\A^1$.
\end{exercise}
\begin{exercise}\label{ex:bline}
List the line bundles on $B\mu_3$, compute their global sections, and determine
which descend to the coarse point. What are their tensor cubes?
\end{exercise}
\begin{exercise}\label{ex:weightcount}
On $\PP(3,1)$, list the monomials in degrees $2,3,4,6$. Compute the
pushforwards of $\OO(2),\OO(3),\OO(-1)$. Which of these stack lines descend?
\end{exercise}
\begin{exercise}\label{ex:fibre}
For $[\A^1_z/\mu_n]\to\A^1_u$, calculate the scheme-theoretic fibre over
$u=0$ and its reduction. Calculate the normal and conormal characters of
the reduced origin substack $B\mu_n$ inside the ambient stack
$[\A^1_z/\mu_n]$.
\end{exercise}
\begin{exercise}\label{ex:generic}
Let $\Gm$ act on $\A^2\setminus0$ with weights $(n,n)$. Show that the
quotient is DM but has stabiliser $\mu_n$ at every geometric point. Why
does this differ from $\PP(n,1)$? Does the coarse space detect the difference?
\end{exercise}
\begin{exercise}\label{ex:push}
In the quadratic-cone example, verify explicitly that the odd module $M$
is generated by $x,y$, has generic rank one, and is not locally free at
the cone vertex. Explain why a stack line with this pushforward cannot
be the pullback of a line bundle on the cone.
\end{exercise}
\begin{exercise}\label{ex:inertia}
The \emph{inertia stack} $I\cX$ parameterises pairs $(\xi,\alpha)$ with
$\alpha\in\Aut(\xi)$. Compute $I[\A^1/\mu_n]$ for scalar multiplication.
Which components come from the nonidentity automorphisms?
\end{exercise}
\begin{exercise}\label{ex:twist}
In the lower-cut toy model replace $d=1$ by a positive integer $d$.
Determine the semistable locus, the stack, its polarisation, and the
criterion for that polarisation to descend. What changes when $n=1$?
\end{exercise}
\begin{exercise}\label{ex:twonormals}
For $n=3$, restrict the two stack normal lines and the two polarisation
lines in the cutting toy model to the common boundary $B\mu_3$.
Write their characters. Which have nonzero invariant vectors?
\end{exercise}

\subsection*{Solutions}
\begin{enumerate}[label=\textbf{\arabic*.},leftmargin=*]
\item For Exercise~\ref{ex:trivial}, the trivial-action stack is
$\A^1\times B\mu_n$ with coarse $\A^1$ and stabiliser $\mu_n$
everywhere. Its restriction to $\Gm$ is $\Gm\times B\mu_n$.
For scalar multiplication, the invariant ring is $k[z^n]$; the coarse
space is again $\A^1$, but only the origin has stabiliser $\mu_n$.
The nonzero open is the ordinary scheme $\Gm$.

\item In Exercise~\ref{ex:torsor}, $z\mapsto z^n$ is a $\mu_n$-torsor on
$\Gm$, so the free-quotient calculation gives the isomorphism.
At zero there are $n$ automorphisms in the quotient stack, whereas a
scheme point has only the identity. Equivalence of stacks preserves
automorphism groups, so the proposed extension is impossible.

\item In Exercise~\ref{ex:bline}, the characters are $0,1,2$ modulo $3$.
Their section spaces are $k,0,0$, and only character $0$ descends.
All three tensor cubes are trivial. A torsion line on a stack can thus
be nontrivial even when its coarse space is a point.

\item In Exercise~\ref{ex:weightcount}, the bases are
\[
 \begin{array}{c|l}
 2&y^2\\
 3&y^3,x\\
 4&y^4,xy\\
 6&y^6,xy^3,x^2.
 \end{array}
\]
Formula~\eqref{eq:weightedpush} gives $\OO$, $\OO(1)$, and
$\OO(-1)$ on the coarse $\PP^1$. Of the three stack lines listed,
only $\OO(3)$ descends. In particular, the invertible pushforward of
$\OO(-1)$ is not proof of descent.

\item In Exercise~\ref{ex:fibre}, the fibre is
$[\Spec k[z]/(z^n)/\mu_n]$ and the reduction is $B\mu_n$.
The normal character is $+1$ and the conormal character is $-1$.
The nilpotent thickening records the multiplicity $n$ in
$\pi^*(0)=n\cD$.

\item In Exercise~\ref{ex:generic}, any nonzero coordinate forces
$t^n=1$, so every stabiliser is $\mu_n$. The coarse space is $\PP^1$,
as follows from the invariant ratio on each chart or regrading the
homogeneous coordinate ring. The difference from $\PP(n,1)$ is generic
isotropy. This quotient is a $\mu_n$-gerbe over $\PP^1$; this description
does not assert a product $\PP^1\times B\mu_n$. A global product
would require a trivialisation of the gerbe.

\item For Exercise~\ref{ex:push}, an odd monomial of positive $x$-exponent
is $x$ times an even monomial; if its $x$-exponent is zero, it is $y$
times an even monomial. Thus $M=Ax+Ay$. Generically,
$y=(xy/x^2)x$, so the rank is one. Modulo the vertex ideal the two
degree-one generators remain independent. Hence the module is not
locally free of rank one there. Moreover its fibre character at the
stack origin is sign, which directly obstructs descent. If it were
pulled back from a coarse line, its pushforward would recover that line
and be locally free, giving a second contradiction.

\item In Exercise~\ref{ex:inertia}, a geometric pair is $(z,\zeta)$ with
$\zeta z=z$. Since $\mu_n$ is abelian, quotienting each fixed locus
does not permute the labels $\zeta$. Thus
\[
 I[\A^1/\mu_n]
 \simeq[\A^1/\mu_n]\ \sqcup\!
 \coprod_{\zeta\ne1}B\mu_n.
\]
The identity component is the original stack; each other automorphism
has fixed locus $0$ and supplies a separate $B\mu_n$ component. This
is a first example of the extra sectors used in orbifold formulas.

\item In Exercise~\ref{ex:twist}, invariance is $na+b=md$. For any
$d>0$, both axes away from the origin have some nonvanishing invariant
positive-degree section, while the origin does not. Hence the
semistable locus and stack remain $\A^2\setminus0$ and $\PP(n,1)$.
The polarisation is now $\OO_{\PP(n,1)}(d)$; it descends exactly when
$n\mid d$. The coarse space remains $\PP^1$. When $n=1$ the action
on $\A^2\setminus0$ is free, the stack is $\PP^1$, and every such
polarisation is the ordinary $\OO_{\PP^1}(d)$.

\item In Exercise~\ref{ex:twonormals}, the lower normal has character
$1$, the upper normal has character $-1=2$ modulo $3$, and both
polarisations have character $1$. None of these four boundary lines
has a nonzero invariant vector. Their cubes are all trivial character
lines and each has a one-dimensional invariant space.
\end{enumerate}

\section{A dictionary and a reading plan}
\label{sec:dictionary}

\begin{longtable}{@{}p{0.29\textwidth}p{0.65\textwidth}@{}}
\toprule
Term & Meaning in these notes\\
\midrule
\endhead
Groupoid & Objects together with invertible arrows; stabilisers are automorphism groups.\\
Stack & Families of objects forming groupoids, with descent for both arrows and objects.\\
Torsor & A family of copies of a group with no preferred identity section.\\
Quotient stack $[X/G]$ & Torsors $P\to T$ with equivariant maps $P\to X$, varying over $T$.\\
Atlas & A representable smooth or étale surjective map from a scheme to the stack.\\
Algebraic stack & A stack with algebraic-space-valued isomorphism sheaves and a smooth atlas.\\
DM stack & An algebraic stack admitting an étale atlas; equivalently its diagonal is unramified.\\
Coarse moduli space & An algebraic space universal for maps to algebraic spaces and preserving geometric isomorphism classes as points.\\
Good moduli space & A quasi-compact map to an algebraic space with exact quasi-coherent pushforward and unchanged structure sheaf. It can identify distinct orbit classes.\\
Inertia & Objects with a chosen automorphism; it records which stabilisers occur in families.\\
Descent of a line bundle & Existence of a coarse line bundle whose pullback is the given stack line. This is stronger than invertibility of its pushforward.\\
Divisor root stack & Families of $n$th roots of a divisor line bundle together with its canonical section.\\
Stack normal line & The normal bundle of the stack boundary, including its stabiliser character.\\
\bottomrule
\end{longtable}

For a second pass, work through the sections in the following order.
\begin{enumerate}
\item Recompute the two affine invariant rings in Section~\ref{sec:quotients}.
Then explain the torsor $s^n=t$ to yourself as a family that cannot be
described by a global root. These are the two independent reasons for
retaining stabilisers and descent.
\item Read the quotient-stack definition alongside the atlas identity
$T\times_{[X/G]}X=P$. This identity explains most of the formal vocabulary.
Use the diagonal-as-isomorphisms description to understand the DM criterion.
\item Work out $\PP(3,1)$ on both charts, including the characters of
$\OO(1)$ and the normal line. Verify the floor in the pushforward formula
and the failure of tensor compatibility.
\item Return to the cut. Keep separate the tensor degree $m$, the group
weight $w$, the stabiliser character modulo its order, and any residual
weight. Derive \eqref{eq:cutinvariance} before using an inequality.
\end{enumerate}

For further reading, Casalaina-Martin--Wise \cite[\S\S2--6]{CMW} develops
the categorical language with many examples. Behrend
\cite[\S\S3.3--3.6]{Behrend} is particularly useful for the quotient,
bundle, and root-stack geometry used here. Vistoli \cite{Vistoli}
provides a more systematic treatment of descent. The Stacks Project
provides precise definitions and theorem statements; its tag references
are more stable than section numbering. Alper \cite{Alper} is the next
step when critical GIT levels and vector-bundle descent become central.
Cadman \cite{Cadman} is the original source for the divisor root construction.

\appendix
\section{Relation to the repository and source conventions}

These notes are a separate introductory document, not proposed manuscript
edits. They use the repository as motivation and as a source of notation.
The mathematical definitions are stated independently, and the examples
are worked out in the text. General existence, local-structure, and
descent theorems have been identified as such rather than claimed to
have elementary proofs here.

The project-specific points of contact are:
\begin{itemize}
\item \path{claude-notes/conventions.tex}, especially its section on
cutting data and names, supplies the convention
$w+\epsilon\ell-md=0$, calligraphic stack notation, and the names of
the two cuts.
\item \path{claude-notes/comparison-theorem.tex}, particularly the
regular/critical-level remarks and the boundary discussion, supplies
the context in which finite boundary stabilisers and polarisation
pushforward matter. The general comparison theorem is not needed for
the elementary toy calculation.
\item \path{SemiprojectiveAlgebraicCutting-integrated.tex} discusses
stack normal lines, finite quotient slices, and possible failure of
coarse polarisation descent. The examples here explain those phenomena
from the beginning.
\end{itemize}

The older integrated manuscript uses a different sign convention for
section weights and the auxiliary direction. The current convention
note records the translation $w_{\mathrm{old}}=-w$ and
$\epsilon_{\mathrm{old}}=-\epsilon$; its old $Y_-$ corresponds to the
current $M^{\le d}$. Thus a cut label or inequality from that older text
must be translated before comparing it with these notes. The invariant
equation in Section~\ref{sec:cutting} is the convention used throughout
this document.

The meaning of ``semiprojective'' also varies in the literature. The
projective-affinisation hypothesis used here is stated separately from
the existence of a contracting action. Nothing in the definition of a
DM stack depends on either property.

\begin{thebibliography}{99}
\bibitem{CMW}
S.~Casalaina-Martin and J.~Wise,
\emph{An introduction to moduli stacks, with a view towards Higgs bundles
on algebraic curves}, arXiv:1708.08124.
\url{https://arxiv.org/abs/1708.08124}.
The saved 147-page PDF in the local Zotero library was consulted;
section and page locators refer to that version.

\bibitem{Behrend}
K.~Behrend, \emph{Introduction to algebraic stacks}, lecture notes,
local Zotero record dated 2014, 127-page version.
\url{https://personal.math.ubc.ca/~behrend/math615A/stacksintro.pdf}.
Relevant material: \S\S3.3--3.6, pp.~110--125. We use the modern coarse
moduli definition given in Section~\ref{sec:coarse}; older conventions
for this terminology differ.

\bibitem{Vistoli}
A.~Vistoli, \emph{Notes on Grothendieck topologies, fibered categories
and descent theory}, arXiv:math/0412512; also in
\emph{Fundamental Algebraic Geometry}, Mathematical Surveys and
Monographs 123, American Mathematical Society, 2005.
\url{https://arxiv.org/abs/math/0412512}.

\bibitem{Stacks}
The Stacks Project Authors, \emph{The Stacks Project}.
\url{https://stacks.math.columbia.edu/}.
Tags used here: 026N (definitions), 04UV (quotient stacks),
04TK (groupoid quotients), 06N3 (DM criterion), and 04UX (coarse-space
literature). Consulted 13 November 2026.

\bibitem{Alper}
J.~Alper, \emph{Good moduli spaces for Artin stacks},
Annales de l'Institut Fourier \textbf{63} (2013), 2349--2402.
\url{https://aif.centre-mersenne.org/articles/10.5802/aif.2833/}.
For vector-bundle descent see Theorem~10.3, journal p.~2386.

\bibitem{Cadman}
C.~Cadman, \emph{Using stacks to impose tangency conditions on curves},
American Journal of Mathematics \textbf{129} (2007), 405--427;
arXiv:math/0312349.
\url{https://arxiv.org/abs/math/0312349}.
For the root construction see \S2.
\end{thebibliography}
\end{document}
